\chapter*{Abstract}
\addcontentsline{toc}{chapter}{Abstract}

This engineering thesis concerns the development of a nonlinear flight dynamics simulation environment for a Boeing 777-like transport aircraft and the design of an autopilot concept based on Nonlinear Model Predictive Control. The aircraft is treated as a large twin-engine wide-body transport configuration inspired by publicly available Boeing 777-300ER data, but the model is not intended to reproduce a certified manufacturer simulation model. The work therefore focuses on engineering consistency, transparent assumptions, and a modular software structure suitable for control design studies.

The thesis describes the adopted coordinate systems, state and input definitions, mass and inertia estimates, standard atmosphere model, first-cut aerodynamic model, propulsion assumptions, trimming requirements, and validation strategy. The implementation is organised as a MATLAB/Simulink-oriented project with reusable data modules, feature tests and reproducible data-processing tools. Particular attention is given to the separation between fixed aircraft properties, such as mass and inertia, and force quantities that depend on environmental assumptions such as gravitational acceleration.

The expected result of the project is a coherent simulation plant that can be trimmed in representative flight conditions and used as a basis for longitudinal and lateral-directional NMPC autopilot modes. The proposed controller architecture is designed to handle input limits, actuator rate limits, state constraints, and trajectory tracking objectives within a finite-horizon optimisation problem.

\textbf{Keywords:} flight dynamics, six-degree-of-freedom model, Boeing 777-like aircraft, nonlinear model predictive control, autopilot, MATLAB, Simulink.
